\section{Lagrangian counts and ordered triple orbits}
\label{sec:orbits}

The contribution of singular differences is nonnegative, so for its upper
bound we may discard the discriminant-sign restrictions after comparing
the signed and unsigned normalizing masses. Their ratio stays bounded
at fixed \(D\). Symmetric matrices then admit a
useful interpretation as Lagrangian subspaces: nullity of a difference becomes
intersection dimension.  We first give the exact counting formulas needed
for this change of model.  We then prove a lower bound for every individual
ordered triple orbit.  This stronger statement is needed for the change of
measure in Section~\ref{sec:planted}.

Throughout this section, dimensions are fixed and $q$ is an odd prime.
The notation $A\asymp_D B$ means that the ratio is bounded above and below
by positive constants depending only on $D$.  In particular, these constants
will be uniform over all the subspaces and form classes that occur.
As before, $b(n)=n(n+1)/2$.

\subsection{Symplectic reduction and elementary counts}

Symplectic bases, quotients, and Lagrangian counts are standard finite
classical geometry; see Taylor~\cite[Theorem~7.4, Lemma~7.5, and Exercises~8.1--8.2]{taylor1992}. We give the arguments
explicitly and then establish the individual ordered-orbit estimate
needed here.

Let $(\mathcal V,\omega)$ be a nondegenerate symplectic space of dimension
$2N$ over $\F_q$.  A subspace is \emph{isotropic} if $\omega$ vanishes on it;
a \emph{Lagrangian} is an isotropic subspace of dimension $N$.
Write $\Lag(\mathcal V)$ for the set of Lagrangians, always equipped with
uniform probability.  Orthogonal complements in this section refer to
$\omega$.  Thus every Lagrangian $L$ satisfies $L^\perp=L$.

We record explicitly the reduction facts used below.  An isotropic basis
$e_1,\ldots,e_t$ extends to a symplectic basis.  To see this, choose vectors
$f_1,\ldots,f_t$ with $\omega(e_i,f_j)=\delta_{ij}$, using nondegeneracy.
If $a_{ij}=\omega(f_i,f_j)$, replace $f_i$ by
$f_i+\sum_k(a_{ki}/2)e_k$.  The new $f_i$ pair to zero with one another,
while their pairings with the $e_i$ are unchanged.  Their joint span is
nondegenerate; repeat in its orthogonal complement.  In particular,
symplectic automorphisms act transitively on isotropic subspaces of any
specified dimension.

For an isotropic $t$-space $S$, the quotient
$S^\perp/S$ is symplectic of dimension $2(N-t)$.  The map
\begin{equation}
 L\longmapsto L/S
 \quad\text{is a bijection from }
 \{L\in\Lag(\mathcal V):S\subseteq L\}
 \text{ to }\Lag(S^\perp/S).
 \label{eq:orbits-containing-reduction}
\end{equation}
More generally, every Lagrangian $L$ has a Lagrangian reduction
\begin{equation}
 \overline L^{\,S}=\big((L\cap S^\perp)+S\big)/S.
 \label{eq:orbits-general-reduction}
\end{equation}
Indeed, if $u=\dim(L\cap S)$, the pairing map $S\to L^*$ has kernel
$S\cap L^\perp=S\cap L$ and rank $t-u$.  Its transpose therefore gives
$\dim(L\cap S^\perp)=N-t+u$.  The space in
Equation~\eqref{eq:orbits-general-reduction} is isotropic and has dimension
$N-t$.  Choosing symplectic coordinates for $S$ and its dual also gives
\begin{equation}
 \mathcal V=(S\oplus S^*)\mathbin{\perp}\mathcal U,
 \qquad \mathcal U\cong S^\perp/S.
 \label{eq:orbits-quotient-extension}
\end{equation}
Here $S^*$ denotes a chosen isotropic dual subspace.  Every symplectic
automorphism of the quotient extends by the identity on $S\oplus S^*$.
In particular, it extends fixing $S$ pointwise.

For $0\le t\le n$, put
\begin{equation}
 \genfrac{[}{]}{0pt}{}{n}{t}_q
 =\prod_{i=0}^{t-1}\frac{q^{n-i}-1}{q^{t-i}-1},
 \qquad
 L_N=\prod_{j=1}^N(q^j+1),\qquad L_0=1.
 \label{eq:orbits-basic-notation}
\end{equation}
An empty product is one.  Counting ordered independent $t$-tuples
in $\F_q^n$ and dividing by the number of ordered bases of a
$t$-space shows that $\genfrac{[}{]}{0pt}{}{n}{t}_q$ counts the
$t$-dimensional subspaces of $\F_q^n$.  We shall also use
\begin{equation}
 |\operatorname{GL}_n(\F_q)|
 =q^{n^2}\prod_{j=1}^n(1-q^{-j}),
 \qquad |\operatorname{GL}_0(\F_q)|=1.
 \label{eq:orbits-gl-order}
\end{equation}

\begin{lemma}[Lagrangian counting formulas]
\label{lem:lagrangian-counts}
In a symplectic $2N$-space, the following formulas hold.
\begin{enumerate}
 \item The number of Lagrangians is $L_N$.
 \item A fixed isotropic $t$-space is contained in exactly $L_{N-t}$
 Lagrangians.  Consequently its containment probability and the number of
 isotropic $t$-spaces are, respectively,
 \begin{equation}
 \vartheta_N(t)=\frac{L_{N-t}}{L_N},\qquad
 I_N(t)=\frac{L_N}{L_{N-t}}\genfrac{[}{]}{0pt}{}{N}{t}_q.
 \label{eq:orbits-containment}
 \end{equation}
 \item For a fixed Lagrangian $A$, the number of Lagrangians $B$ with
 $\dim(A\cap B)=h$ is $\genfrac{[}{]}{0pt}{}{N}{h}_q q^{b(N-h)}$.  Thus
 \begin{equation}
 \pi_{N,h}:=\P\{\dim(A\cap B)=h\}
 =\frac{\genfrac{[}{]}{0pt}{}{N}{h}_q q^{b(N-h)}}{L_N}.
 \label{eq:orbits-pair-probability}
 \end{equation}
 All ordered pairs with the same intersection dimension belong to one
 symplectic orbit.
\end{enumerate}
For fixed $N$, these formulas give
\begin{align}
 L_N&=(1+O_N(q^{-1}))q^{b(N)},
 &\genfrac{[}{]}{0pt}{}{n}{t}_q&=(1+O_n(q^{-1}))q^{t(n-t)},\\
 \vartheta_N(t)&\asymp_N q^{-Nt+\binom t2},
 &I_N(t)&\asymp_N q^{2Nt-(3t^2-t)/2},\\
 \pi_{N,h}&\asymp_N q^{-b(h)}.
 \label{eq:orbits-count-exponents}
\end{align}
\end{lemma}

\begin{proof}
Every line is isotropic.  Count pairs consisting of a line and a Lagrangian
containing it.  Each Lagrangian has $(q^N-1)/(q-1)$ lines, and there are
$(q^{2N}-1)/(q-1)$ lines in $\mathcal V$.  Reduction by a line gives
$L_N(q^N-1)=(q^{2N}-1)L_{N-1}$.  Starting at $L_0=1$ yields the product
in Equation~\eqref{eq:orbits-basic-notation}.
Equation~\eqref{eq:orbits-containing-reduction} gives the containment
count.  Counting pairs $(S,L)$ with $\dim S=t$ and $S\subseteq L$ gives
$I_N(t)L_{N-t}=L_N\genfrac{[}{]}{0pt}{}{N}{t}_q$.

For the pair count, first choose $H=A\cap B$ inside $A$.
In $H^\perp/H$, the image of $B$ must be transverse to the image of $A$,
where \emph{transverse} means having zero intersection.  There are exactly
$q^{b(N-h)}$ such Lagrangians: after choosing complementary coordinate
Lagrangians, they are graphs of symmetric linear maps from one coordinate
space to the other.  Symmetry follows by evaluating the alternating form
on two graph vectors.  This proves Equation~\eqref{eq:orbits-pair-probability}.
To prove pair transitivity, first identify the common subspaces by
symplectic bases.  In their common-space quotient, identify the transverse
pairs by choosing a basis of one member and its dual basis in the other.
The resulting quotient map extends by
Equation~\eqref{eq:orbits-quotient-extension}; containment reduction is a
bijection, so the extension identifies the original pairs.
Finally, factoring the leading powers of $q$ out of the finite products
proves all the displayed estimates, including their endpoint cases.
\end{proof}

\subsection{Comparison with the matrix chart}

Take $N=D=2r$, identify $\mathcal V=\F_q^D\oplus(\F_q^D)^*$, and use
\begin{equation}
 \omega((x,f),(y,g))=g(x)-f(y).
 \label{eq:orbits-coordinate-form}
\end{equation}
For $X\in\mathcal S_D$, write $\Gamma_X=\{(x,Xx):x\in\F_q^D\}$,
where coordinates identify a symmetric matrix with a map to the dual.
These are exactly the Lagrangians transverse to the fixed second
coordinate space.  Denote this chart by $\mathcal A$.  Then
\begin{equation}
 |\mathcal A|=q^{b(D)},\qquad
 \beta_D:=\frac{|\mathcal A|}{L_D}=\prod_{j=1}^D(1+q^{-j})^{-1},
 \qquad
 \dim(\Gamma_X\cap\Gamma_Y)=D-\rank(X-Y).
 \label{eq:orbits-chart}
\end{equation}
Define the unsigned normalized incidence kernel by
\begin{equation}
 J(L,M)=\frac{\1_{\{\dim(L\cap M)=r\}}}{\pi_{D,r}}.
 \label{eq:orbits-unsigned-kernel}
\end{equation}
It has mean one in either variable for every fixed value of the other.
Let
$p_\xi=\P\{\rank X=r,\ \operatorname{sign}(X)=\xi\}$,
so that $K_\xi$ has normalization $p_\xi^{-1}$. This is the quantity
denoted \(\pi_\xi\) in Equation~\eqref{eq:transverse-signed-mass};
the letter \(p\) here distinguishes it from the Lagrangian pair
probability \(\pi_{D,r}\).
Lemma~\ref{lem:symmetric-counts} and
Equation~\eqref{eq:orbits-pair-probability} imply
\begin{equation}
 p_\xi=(1/2+O_D(q^{-1}))q^{-b(r)},\qquad
 \pi_{D,r}=(1+O_D(q^{-1}))q^{-b(r)}.
 \label{eq:orbits-normalizers}
\end{equation}

\begin{lemma}[Chart comparison]
\label{lem:chart-comparison}
Let $B\subseteq\mathscr C$ be a set of active incidences, with prescribed
signs $\xi_{ij}$.  Let $F$ be any nonnegative function of the $35$
Lagrangians indexed by $I\sqcup J$.  For independent uniform matrices
$X_i,Y_j\in\mathcal S_D$ and independent uniform Lagrangians $L_i,M_j$,
\begin{align}
 &\E\left[
 F((\Gamma_{X_i})_i,(\Gamma_{Y_j})_j)
 \prod_{(i,j)\in B}K_{\xi_{ij}}(X_i-Y_j)\right]\notag\\
 &\qquad\le
 \beta_D^{-35}\prod_{(i,j)\in B}\frac{\pi_{D,r}}{p_{\xi_{ij}}}
 \ \E\left[F((L_i)_i,(M_j)_j)
 \prod_{(i,j)\in B}J(L_i,M_j)\right].
 \label{eq:orbits-chart-comparison}
\end{align}
The multiplier is bounded in terms of $D$ alone for all sufficiently large
$q$, uniformly in the active set and its signs.  In particular, this
comparison applies with $F$ the indicator that some pair in $E$ has
positive intersection dimension, or the indicator of any specified
intersection profile.
\end{lemma}

\begin{proof}
On the chart, dropping the sign condition gives the pointwise inequality
\begin{equation}
 K_\xi(X-Y)\le
 \frac{\pi_{D,r}}{p_\xi}J(\Gamma_X,\Gamma_Y).
 \label{eq:orbits-pointwise-comparison}
\end{equation}
The graph variables have exactly the law of the uniform Lagrangian
variables conditioned on all $35$ variables belonging to $\mathcal A$.
That event has probability $\beta_D^{35}$.  Apply
Equation~\eqref{eq:orbits-pointwise-comparison}, express the conditional
expectation by inserting the event's indicator and dividing by its
probability, and then remove the indicator by nonnegativity.  This gives
Equation~\eqref{eq:orbits-chart-comparison} without any boundedness
assumption on $F$ or on the normalized kernels.  Finally
$\beta_D\to1$ and $\pi_{D,r}/p_\xi\to2$.
\end{proof}

\subsection{The cost of an ordered triple orbit}

Let $L_1,L_2,L_3$ be Lagrangians in a symplectic $2D$-space.  Use opposite
indices for their pair intersections and put
\begin{equation}
 \begin{gathered}
 c=\dim(L_1\cap L_2\cap L_3),\qquad
 h_i=\dim(L_j\cap L_k),\\
 s_i=h_i-c,\qquad d=D-c,
 \qquad \{i,j,k\}=\{1,2,3\}.
 \end{gathered}
 \label{eq:orbits-profile}
\end{equation}
These integers form the \emph{dimensional profile}.  Set
\begin{equation}
 C=(D+1)c+\sum_{i=1}^3b(s_i).
 \label{eq:orbits-profile-cost}
\end{equation}
We call a dimensional profile \emph{feasible} if it is realized by an
ordered triple of Lagrangians over the field under consideration.

Here is why the lower bound below concerns individual orbits. Let
\(\mu\) be the independent triple law, and let \(\nu\) be the law
obtained by choosing a uniform center \(Y\) and then, independently
given \(Y\), three uniform Lagrangians meeting \(Y\) in dimension \(r\).
Section~\ref{sec:planted} will verify that the density of \(\nu\)
relative to \(\mu\) is \(\E_Y\prod_i J(L_i,Y)\). This density is
constant on an ordered symplectic orbit \(\mathcal O\), with value
\(\nu(\mathcal O)/\mu(\mathcal O)\). The analogous ratio for a whole
dimensional profile is only the \(\mu\)-weighted average of its orbit
values. A pointwise bound for the center integral therefore needs a
lower bound for \(\mu(\mathcal O)\) for every individual orbit.

\begin{proposition}[Mass of ordered triple orbits]
\label{prop:triple-orbits}
For three independent uniform Lagrangians, every feasible dimensional
profile has probability $O_D(q^{-C})$.  More strongly, there is a constant
$\kappa_D>0$, independent of $q$ and of the particular orbit, such that
for every occurring ordered symplectic orbit $\mathcal O$ of that profile,
\begin{equation}
 \kappa_D q^{-C}\le
 \P\{(L_1,L_2,L_3)\in\mathcal O\}
 \le \kappa_D^{-1}q^{-C}.
 \label{eq:orbits-orbit-mass}
\end{equation}
The constant can be chosen for all odd primes.  Moreover, conditional on
any fixed ordered pair having intersection dimension $h$, the probability
that an independent third Lagrangian completes a specified dimensional
profile is
\begin{equation}
 O_D(q^{-C+b(h)}).
 \label{eq:orbits-conditional-profile}
\end{equation}
\end{proposition}

\begin{proof}
We divide the argument into the common-space count, the classification
above the pair intersections, and the orbit lower bound.

\paragraph{Counting the forced subspaces.}
Choose an isotropic common $c$-space $C_0$ and impose its containment in
all three Lagrangians.  The number of choices times the containment
probability is exactly
\begin{equation}
 I_D(c)\vartheta_D(c)^3\asymp_D q^{-(D+1)c}.
 \label{eq:orbits-common-cost}
\end{equation}
Conditional on these three separate containments, their images in
$C_0^\perp/C_0$ are independent uniform Lagrangians of a symplectic
$2d$-space.  In that quotient we require common intersection zero.
For the next part of the argument denote the reduced Lagrangians again
by $L_i$, and write $H_i=L_j\cap L_k$, of dimension $s_i$.

The $H_i$ are pairwise orthogonal, since any two lie in a common
Lagrangian.  They are also in direct sum.  Indeed, a relation
$x_1+x_2+x_3=0$ with $x_i\in H_i$ puts $x_i$ in $L_i$, since the other
two summands lie there.  Hence each $x_i$ is in the zero common
intersection.  Thus
\begin{equation}
 H_*:=H_1\oplus H_2\oplus H_3\text{ is isotropic},
 \qquad a:=s_1+s_2+s_3\le d,
 \qquad z:=d-a\ge0.
 \label{eq:orbits-pair-sum}
\end{equation}
Choose $H_*$ and then its ordered decomposition into the $H_i$.
The number of such choices is
\begin{equation}
 I_d(a)\frac{|\operatorname{GL}_a(\F_q)|}
 {\prod_i|\operatorname{GL}_{s_i}(\F_q)|}
 \asymp_D
 q^{2da-(3a^2-a)/2+a^2-\sum_i s_i^2}.
 \label{eq:orbits-pair-space-count}
\end{equation}
The stabilizer of an ordered direct-sum decomposition in
$\operatorname{GL}_a$ consists exactly of the block diagonal maps,
which proves the quotient in this formula.
For these choices each $L_i$ must contain
$A_i:=\bigoplus_{j\ne i}H_j$, of dimension $a-s_i$.
Their independent containment probabilities give the additional factor
$\prod_i\vartheta_d(a-s_i)$.  The total exponent is
\begin{align}
 &2da-\frac{3a^2-a}{2}+a^2-\sum_i s_i^2
 +\sum_i\left[-d(a-s_i)+\binom{a-s_i}{2}\right]
 =-\sum_i b(s_i).
 \label{eq:orbits-pair-cost-cancellation}
\end{align}
Multiplying with Equation~\eqref{eq:orbits-common-cost}, and ignoring
all conditions beyond these containments, proves the asserted profile
upper bound.

The profile upper bound is now established. The center-factor comparison
also needs a lower bound for each individual ordered orbit; the remaining
argument supplies that bound while keeping the intersections exact.

\paragraph{Splitting off the pair intersections.}
We next identify all invariants of an actual common-zero triple.  Fix
$i$.  The pairing $H_i\to L_i^*$ is injective, since its kernel is
$H_i\cap L_i^\perp=H_i\cap L_i=0$.  Its transpose is therefore
surjective.  Choose an $s_i$-space $Q_i\subseteq L_i$ paired perfectly
with $H_i$.  Both spaces are isotropic, so
$B_i=H_i\oplus Q_i$ is a nondegenerate symplectic $2s_i$-space.
Each $L_\ell$ contains a Lagrangian of $B_i$: it contains $Q_i$ when
$\ell=i$ and $H_i$ otherwise.  Call that subspace $K_\ell$.
Under the orthogonal decomposition $\mathcal V=B_i\mathbin\perp B_i^\perp$,
the projection of a vector $x\in L_\ell$ to $B_i$ is orthogonal to
$K_\ell$, hence belongs to $K_\ell$.  As $K_\ell\subseteq L_\ell$,
subtracting the projection gives the actual splitting
\begin{equation}
 L_\ell=K_\ell\oplus(L_\ell\cap B_i^\perp),
 \qquad \ell=1,2,3.
 \label{eq:orbits-forced-splitting}
\end{equation}
For $j\ne i$, the space $H_j$ lies in $B_i^\perp$: it is orthogonal to
$H_i$, and is orthogonal to $Q_i$ because $H_j,Q_i\subseteq L_i$.
Thus the other pair intersections survive in the orthogonal complement,
and the construction can be repeated there.  The zero-dimensional
blocks require no choice.

After splitting all three blocks, the remaining ordered triple is
transverse in a symplectic space of dimension $2z$.  The space $H_*$
is Lagrangian in the sum of the three blocks, so this remaining triple
is precisely the reduction of the original triple to $H_*^\perp/H_*$
as in Equation~\eqref{eq:orbits-general-reduction}.
Symplectic bases identify any two corresponding pair blocks, respecting
all three members of the ordered triple.  A symplectic map identifying
the residual ordered triples extends over these block maps.  Consequently,
the dimensions $s_i$ and the residual ordered orbit determine the entire
common-zero ordered orbit.  For $c>0$, first identify the common spaces;
Equation~\eqref{eq:orbits-containing-reduction} and the extension in
Equation~\eqref{eq:orbits-quotient-extension} show that the same assertion
lifts to the original space.  In particular there is no further invariant
arising from extensions between the blocks.

\paragraph{The mass of a residual orbit.}
An ordered transverse pair of Lagrangians in dimension $2z$ can be
identified with the two coordinate spaces.  A third Lagrangian
transverse to both is then the graph of a nonsingular symmetric form.
The stabilizer of the ordered pair consists of the maps induced by
$g\in\operatorname{GL}_z(\F_q)$ on one coordinate space and its inverse
dual on the other; its action on forms is congruence.  Thus residual
ordered orbits correspond exactly to congruence orbits
$\mathcal O_{\mathrm{form}}$ of nonsingular symmetric forms.
For this classification of ordered transverse triples, see also
Kramer and Tent~\cite[Theorem~6 and Section~2.3]{kramertent2010}.
For three independent uniform residual Lagrangians, a prescribed
residual orbit has probability exactly
\begin{equation}
 \frac{q^{b(z)}|\mathcal O_{\mathrm{form}}|}{L_z^2}.
 \label{eq:orbits-residual-mass}
\end{equation}
Indeed, after the first Lagrangian is chosen, the second is transverse
with probability $q^{b(z)}/L_z$; for each such ordered pair there are
$|\mathcal O_{\mathrm{form}}|$ choices of the third among $L_z$
Lagrangians.

This probability has a positive lower bound depending only on $z$,
uniformly in the form class.  Diagonalize a nonsingular symmetric form $F$
and choose an orthogonal basis with specified nonzero norms.  An
isometry is determined by the images of this basis.  At the stage when
the remaining nondegenerate orthogonal complement has dimension $k$,
there are at most $2q^{k-1}$ vectors of a specified nonzero norm: after
diagonalizing, fix $k-1$ coordinates and solve a quadratic equation in
the last one.  Each chosen vector has nonzero norm, so the next
orthogonal complement is again nondegenerate.  Hence
\begin{equation}
 |\operatorname{Isom}(F)|\le
 2^z q^{z(z-1)/2},\qquad
 |\mathcal O_{\mathrm{form}}|
 \ge 2^{-z}q^{b(z)}\prod_{j=1}^z(1-q^{-j}).
 \label{eq:orbits-isometry-bound}
\end{equation}
Orbit--stabilizer and Equation~\eqref{eq:orbits-gl-order} give the second
inequality.  For example, for every odd prime the lower bound in
Equation~\eqref{eq:orbits-residual-mass} is at least
\begin{equation}
 \gamma_z:=2^{-z}
 \frac{\prod_{j=1}^z(1-3^{-j})}
 {\prod_{j=1}^z(1+3^{-j})^2}>0.
 \label{eq:orbits-residual-constant}
\end{equation}
For $z=0$, there is one residual Lagrangian and one empty form;
$L_0=1$, $|\mathcal O_{\mathrm{form}}|=1$, and
Equation~\eqref{eq:orbits-residual-mass} is exactly one.  We take
$\gamma_0=1$.

We have identified the residual form class and bounded its probability
in its own symplectic space. We must now lift that estimate through
the chosen pair-intersection spaces while keeping their intersections
exact. We condition each Lagrangian separately before imposing the
residual orbit.

\paragraph{Independent residual laws and exact intersections.}
Fix candidate spaces $H_i$ as in
Equation~\eqref{eq:orbits-pair-space-count}, and sample the three
Lagrangians independently subject to $A_i\subseteq L_i$.
Require, separately for each sample, the avoidance event
\begin{equation}
 \mathcal E_i=\{L_i\cap H_*=A_i\}.
 \label{eq:orbits-avoidance}
\end{equation}
In the containment quotient $A_i^\perp/A_i$, of half-dimension
$d-(a-s_i)=z+s_i$, failure means meeting the fixed isotropic
$s_i$-space $H_*/A_i$ in a line.  If $s_i>0$, a union bound and
Equation~\eqref{eq:orbits-containment} give
\begin{equation}
 \P(\mathcal E_i^c\mid A_i\subseteq L_i)
 \le \frac{q^{s_i}-1}{q-1}\frac1{q^{z+s_i}+1}
 <\frac{q^{-z}}{q-1}\le\frac12.
 \label{eq:orbits-avoidance-probability}
\end{equation}
If $s_i=0$, the failure probability is zero.  In particular, each
avoidance probability is at least $1/2$ and is $1-O_D(q^{-1})$.
The containments and avoidance events concern each sample separately,
so after conditioning on all of them the three samples remain
independent.

Let $\widehat L_i$ be the reduction of $L_i$ to $H_*^\perp/H_*$.
Each conditional $\widehat L_i$ is uniform on the Lagrangians of this
quotient.  To verify this, extend any symplectic automorphism of the
quotient by the identity on $H_*$ and a chosen dual, as in
Equation~\eqref{eq:orbits-quotient-extension}.  The resulting map fixes
every $H_i$ and $A_i$, preserves containment and avoidance, and takes
reductions to their images under the prescribed quotient map.  The
conditional law is therefore invariant under the full symplectic
group of the quotient.  That group is transitive on its Lagrangians,
so the law is uniform.  Since the reduction is a function of each
sample individually, the three reductions are also independent.

Require their ordered residual orbit to be the one of the fixed target
orbit $\mathcal O$.  By Equation~\eqref{eq:orbits-residual-mass}, this
has conditional probability at least $\gamma_z$.  These conditions
force the selected pair intersections to be exact.  For example, if
$x\in L_i\cap L_j$, then $x$ is orthogonal to every $H_\ell$:
for each $\ell$, at least one of $L_i,L_j$ contains $H_\ell$.
Thus $x\in H_*^\perp$, and its image lies in
$\widehat L_i\cap\widehat L_j=0$.  It follows that $x\in H_*$.
Avoidance now gives
\begin{equation}
 L_i\cap L_j\subseteq A_i\cap A_j=H_k,
 \qquad \{i,j,k\}=\{1,2,3\}.
 \label{eq:orbits-exact-pairs}
\end{equation}
The reverse inclusion was imposed from the start.  Hence the pair
intersections are exactly the $H_k$, and the common intersection is
zero.  The splitting argument proves that every resulting triple
belongs to the prescribed ordered orbit in the common-space quotient,
and then to $\mathcal O$ in the original space.

For each choice of $C_0,H_1,H_2,H_3$, the required additional probability
is at least $2^{-3}\gamma_z$.  Successful choices cannot count a triple
twice: its exact common space is $C_0$ and its exact pair spaces in the
quotient are the $H_i$.  We have proved the explicit lower bound
\begin{align}
 \P(\mathcal O)\ge{}&2^{-3}\gamma_z\,
 I_D(c)\vartheta_D(c)^3
 I_d(a)\frac{|\operatorname{GL}_a(\F_q)|}
 {\prod_i|\operatorname{GL}_{s_i}(\F_q)|}
 \prod_i\vartheta_d(a-s_i).
 \label{eq:orbits-explicit-lower-bound}
\end{align}
Equations~\eqref{eq:orbits-common-cost} and
\eqref{eq:orbits-pair-cost-cancellation} now give the lower bound in
Equation~\eqref{eq:orbits-orbit-mass}.  The upper bound follows from
the profile upper bound already proved.  All product bounds are
uniform for $q\ge3$ and involve dimensions at most $D$; the explicit
$\gamma_z$ is uniform over actual form classes.  Taking minima and
maxima over these finitely many dimension tuples gives a single
constant $\kappa_D$.

Finally, pair transitivity in Lemma~\ref{lem:lagrangian-counts} makes
the probability of completing a given profile the same for every
fixed ordered pair of dimension $h$.  It equals the unconditional
profile probability divided by $\pi_{D,h}$ when the prescribed pair
dimension is $h$, and is zero otherwise.  The profile upper bound
and Equation~\eqref{eq:orbits-pair-probability} give
Equation~\eqref{eq:orbits-conditional-profile}.
\end{proof}
